\documentclass{article}
\usepackage[utf8]{inputenc}
\usepackage[T1]{fontenc}
\usepackage{lmodern}
\usepackage[polish]{babel}
\usepackage{amsmath}
\usepackage{tikz}
\usepackage{algorithm}
\usepackage{algpseudocode}
\usepackage{hyperref}
\usepackage{float}
\usepackage{graphicx}
\usepackage{mathtools}
\usepackage{amsmath}
\usepackage{amsfonts}
\usepackage{amsmath}
\usepackage{amsmath}
\usepackage{booktabs}
\usepackage[margin=1in]{geometry}

\title{Języki Formalne i Techniki Translacji}
\author{Wojciech Typer 279730}
\date{}

\begin{document}
\maketitle

\section*{Zadanie 3 lista 6}

\subsection*{1. Eliminacja rekurencji lewostronnej}
Oryginalna gramatyka posiada rekurencję lewostronną w produkcjach dla $E$ oraz $T$. Po zastosowaniu standardowego algorytmu eliminacji ($A \to A\alpha \mid \beta \implies A \to \beta A', A' \to \alpha A' \mid \epsilon$), otrzymujemy gramatykę $G'$:

\begin{itemize}
    \item $E \to T E'$
    \item $E' \to \text{or } T E' \mid \epsilon$
    \item $T \to F T'$
    \item $T' \to \text{and } F T' \mid \epsilon$
    \item $F \to \text{not } F \mid ( E ) \mid \text{true} \mid \text{false}$
\end{itemize}

\subsection*{2. Lewostronna faktoryzacja}
Lewostronna faktoryzacja jest stosowana, gdy dla jednego nieterminala istnieją co najmniej dwie produkcje zaczynające się od tego samego symbolu. W gramatyce $G'$:
\begin{itemize}
    \item Produkcje dla $E, E', T, T'$ zaczynają się od unikalnych symboli (włączając $\epsilon$).
    \item Produkcje dla $F$ zaczynają się odpowiednio od: \texttt{not}, \texttt{(}, \texttt{true}, \texttt{false}. 
\end{itemize}
Zatem \textbf{lewostronna faktoryzacja nie jest wymagana}, ponieważ nie isteieje wiele produkcji dla jednego nieterminala zaczynających się od tego samego symbolu.

\subsection*{3. Zbiory FIRST i FOLLOW}
Wyznaczone zbiory dla przekształconej gramatyki (gdzie $\epsilon$ oznacza słowo puste, a \$ koniec ciągu):

\begin{table}[H]
\centering
\begin{tabular}{@{}lll@{}}
\toprule
\textbf{Symbol} & \textbf{FIRST} & \textbf{FOLLOW} \\ \midrule
$E$  & \{ \texttt{not, (, true, false} \} & \{ \texttt{\$, )} \} \\
$E'$ & \{ \texttt{or}, $\epsilon$ \}         & \{ \texttt{\$, )} \} \\
$T$  & \{ \texttt{not, (, true, false} \} & \{ \texttt{or, \$, )} \} \\
$T'$ & \{ \texttt{and}, $\epsilon$ \}        & \{ \texttt{or, \$, )} \} \\
$F$  & \{ \texttt{not, (, true, false} \} & \{ \texttt{and, or, \$, )} \} \\ \bottomrule
\end{tabular}
\end{table}

\subsection*{4. Czy gramatyka jest LL(1)?}
Tak, otrzymana gramatyka \textbf{jest typu LL(1)}. Dla każdej produkcji $A \to \alpha \mid \beta$:
\begin{enumerate}
    \item $FIRST(\alpha) \cap FIRST(\beta) = \emptyset$.
    \item $FIRST(\alpha) \cap FOLLOW(A) = \emptyset$.
\end{enumerate}

\begin{itemize}
    \item \textbf{Nieterminal $E'$ ($E' \to \text{or } T E' \mid \varepsilon$):} \\
    $\alpha = \text{or } T E' \implies FIRST(\alpha) = \{\text{or}\}$ \\
    $\beta = \varepsilon \implies FIRST(\beta) = \{\varepsilon\}$ \\
    $FOLLOW(E') = \{\$, )\}$ \\
    $\{\text{or}\} \cap \{\$, )\} = \emptyset$. Warunek spełniony.

    \item \textbf{Nieterminal $T'$ ($T' \to \text{and } F T' \mid \varepsilon$):} \\
    $\alpha = \text{and } F T' \implies FIRST(\alpha) = \{\text{and}\}$ \\
    $\beta = \varepsilon \implies FIRST(\beta) = \{\varepsilon\}$ \\
    $FOLLOW(T') = \{\text{or}, \$, )\}$ \\
    $\{\text{and}\} \cap \{\text{or}, \$, )\} = \emptyset$. Warunek spełniony.

    \item \textbf{Nieterminal $F$ ($F \to \text{not } F \mid (E) \mid \text{true} \mid \text{false}$):} \\
    Wszystkie alternatywy zaczynają się od unikalnych terminali: $\{\text{not}\}, \{(\}, \{\text{true}\}, \{\text{false}\}$, więc ich zbiory $FIRST$ są rozłączne.
\end{itemize}

Te warunki są spełnione, co oznacza, że nasza gramatyka jest LL(1)

\subsection*{5. Tablica parsera przewidującego}

\begin{table}[H]
\centering
\small
\begin{tabular}{|c|c|c|c|c|c|c|c|c|}
\hline
 & \textbf{or} & \textbf{and} & \textbf{not} & \textbf{(} & \textbf{)} & \textbf{true} & \textbf{false} & \textbf{\$} \\ \hline
\textbf{E} & & & $E \to T E'$ & $E \to T E'$ & & $E \to T E'$ & $E \to T E'$ & \\ \hline
\textbf{E'} & $E' \to \text{or } T E'$ & & & & $E' \to \epsilon$ & & & $E' \to \epsilon$ \\ \hline
\textbf{T} & & & $T \to F T'$ & $T \to F T'$ & & $T \to F T'$ & $T \to F T'$ & \\ \hline
\textbf{T'} & $T' \to \epsilon$ & $T' \to \text{and } F T'$ & & & $T' \to \epsilon$ & & & $T' \to \epsilon$ \\ \hline
\textbf{F} & & & $F \to \text{not } F$ & $F \to (E)$ & & $F \to \text{true}$ & $F \to \text{false}$ & \\ \hline
\end{tabular}
\end{table}

\end{document}